\chapter{Can offline routing reduce the congestion?}
\label{ch:routing}

Experiment~2 (Chapter~\ref{ch:heatmap}) showed that at the modelled capacity of four memories nearly all overload is relay overload: it is caused by routes passing through the
central cores, not by routes starting or ending there. That raises a question about the routes. The mapper chooses, for every qubit move, one Dijkstra shortest path (fewest hops, link
reliability as the tie-break), without looking at what the other moves of the layer are doing. Many alternative shortest paths exist: between opposite corners of the $4\times4$ mesh there are
$\binom{6}{3}=20$. Could choosing them \emph{offline}, with the mapping untouched, spread the relay load and reduce the congestion?

\section{Method}
The mapping is held fixed: every qubit still moves from the same core to the same core in the same layer. Only the path of each move (a \emph{route}) is re-chosen, layer by layer
(layers are independent because a layer's routes start and end at fixed cores). For each layer that has any overload, the choice of one candidate path per route is made by solving a small
mixed-integer program exactly, minimising the \emph{excess}, the demand above capacity summed over cores. The candidates are all simple paths of at most the shortest length plus $0$ or $2$ hops,
and the mapper's own path is always a candidate, so the result is never worse than the baseline.

Two lower bounds hold for \emph{any} routing, so they show how much is left to gain. Endpoints do not depend on the path, so the demand of a core is at least its endpoint demand; and a route of shortest
length $k$ puts at least $2k$ hop touches ($k+1$ route touches) on the cores, so the total excess is at least the minimum total demand minus (cores $\times$ capacity). Both are tested.

Demand is counted two ways: by hops (the archive's counting) and by routes, the number of qubits in flight through a core, which is what the per-qubit scheduler experiences.

\section{Structural result}
\begin{table}[H]
\centering\small
\begin{tabular}{@{}rrrrrrr@{}}
\toprule
Capacity & Extra hops & Excess now & Excess after re-routing & Lower bound & Overloaded layers & Added hops \\
\midrule
1 & 0 & 27043 & 21894 (19.0\%) & 18336 & 4110 $\to$ 3801 & 0.00\% \\
1 & 2 & 27043 & 21018 (22.3\%) & 18336 & 4110 $\to$ 3621 & 2.55\% \\
2 & 0 & 10385 & 6771 (34.8\%) & 4995 & 1641 $\to$ 978 & 0.00\% \\
2 & 2 & 10385 & 6288 (39.5\%) & 4995 & 1641 $\to$ 840 & 1.48\% \\
3 & 0 & 4774 & 1791 (62.5\%) & 846 & 855 $\to$ 420 & 0.00\% \\
3 & 2 & 4774 & 1497 (68.6\%) & 846 & 855 $\to$ 327 & 0.98\% \\
4 & 0 & 2216 & 348 (84.3\%) & 78 & 482 $\to$ 102 & 0.00\% \\
4 & 2 & 2216 & 237 (89.3\%) & 78 & 482 $\to$ 48 & 0.36\% \\
\bottomrule
\end{tabular}
\caption{Offline re-routing, mesh, route counting, pooled over ten benchmarks and three mapping seeds. Excess is demand above capacity.}
\label{tab:routing-mesh}
\end{table}

Pooled over the ten benchmarks, the best routing removes most of the excess at the modelled capacity: with route counting at capacity four, shortest paths alone cut the excess from 2\,216 to 348 (an 84.3\% reduction)
and the layers with any overload from 482 to 102; allowing two extra hops cuts it to 237 (89.3\%) at a cost of 0.36\% more hops. The gain shrinks as the capacity shrinks (19.0\% at capacity one),
because there the total relay volume, not the choice of path, is the limit, which the volume bound shows. With hop counting the reductions are smaller (65.6\% at capacity four).

\section{Simulated effect}
The structural gain would be an artefact if it did not survive the simulator, so the re-routed traces were run in ACE next to the original routes, on-demand, at the budget they were optimised for, with the same
mapping, endpoints, physics seeds and parameters. Three mapping seeds and three physics seeds give nine paired runs per cell.
\begin{table}[H]
\centering\small
\resizebox{\ifdim\width>\linewidth\linewidth\else\width\fi}{!}{%
\begin{tabular}{@{}lrrrrr@{}}
\toprule
Workload & Mem. & Original & Re-routed & +2 hops & Waited \\
\midrule
Graph state & 2 & 2.72 & 1.53 (+42.3\% $\pm$9.9; 9/9) & 1.41 (+46.6\% $\pm$8.9; 9/9) & 0.281 $\to$ 0.086 \\
Graph state & 3 & 1.53 & 1.15 (+22.9\% $\pm$11.8; 9/9) & 1.19 (+19.6\% $\pm$13.5; 9/9) & 0.121 $\to$ 0.000 \\
Graph state & 4 & 1.24 & 1.16 (+6.0\% $\pm$7.6; 7/9) & 1.21 (+1.9\% $\pm$8.9; 4/9) & 0.029 $\to$ 0.000 \\
Draper QFT adder & 2 & 2.16 & 1.69 (+21.1\% $\pm$7.9; 9/9) & 1.48 (+31.1\% $\pm$6.8; 9/9) & 0.211 $\to$ 0.115 \\
Draper QFT adder & 3 & 1.30 & 1.16 (+10.3\% $\pm$3.3; 9/9) & 1.16 (+10.5\% $\pm$3.7; 9/9) & 0.055 $\to$ 0.007 \\
Draper QFT adder & 4 & 1.15 & 1.13 (+1.3\% $\pm$1.3; 8/9) & 1.14 (+1.1\% $\pm$0.8; 9/9) & 0.008 $\to$ 0.000 \\
QAOA & 2 & 3.42 & 2.24 (+34.5\% $\pm$2.4; 9/9) & 2.12 (+38.0\% $\pm$2.4; 9/9) & 0.376 $\to$ 0.206 \\
QAOA & 3 & 1.74 & 1.25 (+28.4\% $\pm$2.3; 9/9) & 1.19 (+31.4\% $\pm$2.2; 9/9) & 0.172 $\to$ 0.037 \\
QAOA & 4 & 1.27 & 1.13 (+10.9\% $\pm$1.7; 9/9) & 1.13 (+11.5\% $\pm$1.6; 9/9) & 0.058 $\to$ 0.003 \\
QFT & 2 & 4.72 & 3.49 (+26.1\% $\pm$2.4; 9/9) & 3.03 (+35.9\% $\pm$2.1; 9/9) & 0.430 $\to$ 0.312 \\
QFT & 3 & 2.39 & 1.53 (+35.9\% $\pm$2.2; 9/9) & 1.47 (+38.2\% $\pm$2.3; 9/9) & 0.249 $\to$ 0.097 \\
QFT & 4 & 1.59 & 1.19 (+25.2\% $\pm$2.4; 9/9) & 1.17 (+26.6\% $\pm$2.0; 9/9) & 0.130 $\to$ 0.018 \\
Random circuit & 2 & 1.27 & 1.13 (+10.6\% $\pm$1.3; 9/9) & 1.12 (+11.7\% $\pm$1.3; 9/9) & 0.042 $\to$ 0.007 \\
Random circuit & 3 & 1.12 & 1.11 (+1.2\% $\pm$0.3; 9/9) & 1.11 (+1.3\% $\pm$0.3; 9/9) & 0.006 $\to$ 0.000 \\
Random circuit & 4 & 1.11 & 1.11 (+0.2\% $\pm$0.1; 9/9) & 1.11 (+0.2\% $\pm$0.2; 8/9) & 0.001 $\to$ 0.000 \\
\bottomrule
\end{tabular}}
\caption{Simulated on-demand latency (ms), mesh: original routes against offline re-routing. Brackets: paired reduction in latency (positive = faster), sd, pairs faster of nine.}
\label{tab:routing-sim-mesh}
\end{table}

Re-routing alone makes on-demand requests faster in nearly every pair. At two memories the latency falls by 21\% to 42\% on the four congested workloads (graph state, Draper, QAOA, QFT) and by 11\% on the random circuit; at the modelled four memories by 25\% for QFT, 11\% for QAOA and 6\% for graph state,
and by 1\% or less where there was little congestion. The fraction of requests that had to wait falls accordingly. Two extra hops help at two memories (QFT 26\% to 36\%) and change little at four (within about 1.5 percentage points, except graph state, where they cost 4 points): the added hops cost about as much as the avoided waiting saves.

\section{The torus}
On the torus, which has more paths between cores, the optimum is closer to the lower bound still: at capacity four the excess falls from 493 to 15 (97.0\%), against a lower bound of 6, and with two extra hops it reaches the bound.
\begin{table}[H]
\centering\small
\begin{tabular}{@{}rrrrrrr@{}}
\toprule
Capacity & Extra hops & Excess now & Excess after re-routing & Lower bound & Overloaded layers & Added hops \\
\midrule
1 & 0 & 20202 & 15858 (21.5\%) & 14847 & 3826 $\to$ 3339 & 0.00\% \\
1 & 2 & 20202 & 15387 (23.8\%) & 14847 & 3826 $\to$ 3312 & 1.73\% \\
2 & 0 & 6108 & 3351 (45.1\%) & 2847 & 1319 $\to$ 702 & 0.00\% \\
2 & 2 & 6108 & 3108 (49.1\%) & 2847 & 1319 $\to$ 630 & 0.93\% \\
3 & 0 & 1852 & 363 (80.4\%) & 303 & 562 $\to$ 162 & 0.00\% \\
3 & 2 & 1852 & 342 (81.5\%) & 303 & 562 $\to$ 159 & 0.09\% \\
4 & 0 & 493 & 15 (97.0\%) & 6 & 220 $\to$ 15 & 0.00\% \\
4 & 2 & 493 & 6 (98.8\%) & 6 & 220 $\to$ 6 & 0.04\% \\
\bottomrule
\end{tabular}
\caption{Offline re-routing, torus, route counting, pooled over ten benchmarks and three mapping seeds. Excess is demand above capacity.}
\label{tab:routing-torus}
\end{table}

Simulated on the torus, re-routing makes requests faster at two memories (latency falls by 15\% for QFT, 32\% for QAOA, 36\% for Draper) but little at the modelled four (by 6.8\% for QFT, 3.2\% for QAOA, about 2\% or nothing for the others), because the torus is already close to its latency floor there.
\begin{table}[H]
\centering\small
\resizebox{\ifdim\width>\linewidth\linewidth\else\width\fi}{!}{%
\begin{tabular}{@{}lrrrrr@{}}
\toprule
Workload & Mem. & Original & Re-routed & +2 hops & Waited \\
\midrule
Graph state & 2 & 1.86 & 1.27 (+31.2\% $\pm$7.2; 9/9) & 1.28 (+31.0\% $\pm$6.9; 9/9) & 0.157 $\to$ 0.033 \\
Graph state & 3 & 1.25 & 1.10 (+11.0\% $\pm$7.4; 8/9) & 1.12 (+10.3\% $\pm$8.7; 8/9) & 0.033 $\to$ 0.000 \\
Graph state & 4 & 1.11 & 1.12 (-0.4\% $\pm$1.7; 1/9) & 1.11 (+0.3\% $\pm$3.9; 2/9) & 0.000 $\to$ 0.000 \\
Draper QFT adder & 2 & 2.25 & 1.45 (+35.6\% $\pm$1.9; 9/9) & 1.26 (+44.0\% $\pm$1.4; 9/9) & 0.226 $\to$ 0.075 \\
Draper QFT adder & 3 & 1.28 & 1.15 (+10.5\% $\pm$1.8; 9/9) & 1.14 (+10.7\% $\pm$1.0; 9/9) & 0.053 $\to$ 0.007 \\
Draper QFT adder & 4 & 1.13 & 1.11 (+1.8\% $\pm$0.6; 9/9) & 1.11 (+1.6\% $\pm$0.8; 9/9) & 0.008 $\to$ 0.000 \\
QAOA & 2 & 2.47 & 1.66 (+32.3\% $\pm$4.7; 9/9) & 1.62 (+34.2\% $\pm$4.7; 9/9) & 0.265 $\to$ 0.111 \\
QAOA & 3 & 1.40 & 1.15 (+17.6\% $\pm$4.6; 9/9) & 1.14 (+18.3\% $\pm$4.3; 9/9) & 0.083 $\to$ 0.004 \\
QAOA & 4 & 1.17 & 1.13 (+3.2\% $\pm$2.8; 9/9) & 1.13 (+3.1\% $\pm$2.6; 9/9) & 0.016 $\to$ 0.000 \\
QFT & 2 & 2.79 & 2.37 (+14.7\% $\pm$7.3; 9/9) & 2.42 (+12.7\% $\pm$7.8; 9/9) & 0.303 $\to$ 0.221 \\
QFT & 3 & 1.51 & 1.21 (+19.3\% $\pm$6.0; 9/9) & 1.24 (+17.3\% $\pm$6.1; 9/9) & 0.102 $\to$ 0.023 \\
QFT & 4 & 1.21 & 1.13 (+6.8\% $\pm$3.2; 9/9) & 1.12 (+7.3\% $\pm$3.0; 9/9) & 0.028 $\to$ 0.001 \\
Random circuit & 2 & 1.19 & 1.11 (+6.1\% $\pm$0.9; 9/9) & 1.11 (+6.3\% $\pm$0.9; 9/9) & 0.022 $\to$ 0.002 \\
Random circuit & 3 & 1.11 & 1.11 (+0.4\% $\pm$0.1; 9/9) & 1.11 (+0.3\% $\pm$0.1; 9/9) & 0.001 $\to$ 0.000 \\
Random circuit & 4 & 1.11 & 1.11 (-0.0\% $\pm$0.1; 1/9) & 1.11 (+0.0\% $\pm$0.1; 3/9) & 0.000 $\to$ 0.000 \\
\bottomrule
\end{tabular}}
\caption{Simulated on-demand latency (ms), torus: original routes against offline re-routing. Brackets: paired reduction, sd, pairs faster of nine.}
\label{tab:routing-sim-torus}
\end{table}


\section{What can and cannot be concluded}
Offline routing \emph{can} reduce congestion: with the mapping fixed, an exact offline choice of paths removes most of the excess demand at the modelled capacity and the simulated latencies fall with it. These
are results for an offline optimum with global knowledge of every layer, computed per layer for the capacity the hardware has; a practical routing algorithm would be a heuristic for the same problem and would be bounded by it.
Not shown here: the effect on pre-generation policies (only on-demand was simulated), joint optimisation of mapping and routing, fidelity effects (the on-demand fidelity does not depend on waiting), and a routing policy that does not know
the layer in advance. Whether to pursue the research further is therefore a question about those directions, and the bounds above say how much headroom each could have.
